\chapter*{The decision before the model}
\addcontentsline{toc}{chapter}{The decision before the model}

A bank considering a higher deposit rate faces a deceptively simple question:
will the extra funding be worth the interest it must pay? Customers who choose
a generous rate may already intend to move money. An increase in recorded
balances may be a transfer from another account at the same bank. Competitors
may change their prices at the same time. A forecast of balances, however
accurate, cannot by itself separate these explanations.

Causal analysis asks how the bank's choice changes an outcome relative to a
specified alternative. The alternative matters as much as the proposed action.
An invitation to open an account, activation of that account, and a change in
the rate paid on an existing balance are different interventions. Their costs,
beneficiaries and effects need not coincide. A useful analysis begins with this
economic distinction and carries it through the data, statistical assumptions
and financial calculation.

The same discipline applies beyond deposits. A collections team choosing
when to contact a borrower must account for information revealed by earlier
contacts. A lender changing an eligibility threshold must distinguish the
effect on marginal applicants from the effect on all customers. Treasury must
separate an external funding shock from the bank's own response. In each case,
the bank wants to compare feasible decisions, and the available observations
may support only part of that comparison.

The argument develops from these choices. Economic questions lead to causal
models; the models reveal which comparisons the data can identify; estimation
then measures those comparisons and their uncertainty. Repeated decisions and
imperfect observations complicate the analysis. Existing software offers many
of the necessary methods, but its assumptions and definitions must match the
question. Language models may help economists formulate questions and examine
evidence. Public data and controlled numerical studies can establish which
ideas merit a prototype before a bank pays for a larger experiment.

The proposal is for a common analytical service that makes this work easier to
repeat without changing the underlying economics. Its case rests on the
quality and cost of completed studies, including studies that show why an
effect cannot be learned from the available information. The reference
alternative is a capable quantitative team using established methods. A new
service must justify its additional cost against that alternative.

The mathematical development remains part of the argument. Symbols identify
economic quantities, assumptions state what the comparison requires, and
examples expose where a conclusion fails. Constructed examples and simulated
populations are described as such. Findings from public data retain their
population and measurement limits. Proposed capabilities are distinguished
from the calculations that have actually been performed. The literature
cutoff is 7 October 2026; source editions and acquisition dates are documented
separately.

\chapter*{The argument in nine parts}
\addcontentsline{toc}{chapter}{The argument in nine parts}

Part I starts with a bank's decision and develops the distinction between
selection, identification and estimation. Part II follows customers and
institutions through time. Part III considers uncertain causal relationships,
imperfect records and the transfer of results between populations.

Parts IV and V compare available methods and software. Discovery and
identification come first, followed by effect estimation, repeated decisions,
panel comparisons and commercial choices. Detailed package discussions sit
beside the questions they address, so that a reader can compare alternatives
without following the order in which they were acquired.

Part VI examines assistance from language models: turning a business question
into a research design, drawing and questioning causal diagrams, supporting
discovery, and explaining results. Part VII develops practical use through
research libraries, command-line tools and the Model Context Protocol. These
implementation choices follow the economic and statistical requirements.

Part VIII brings public data and banking studies together. It asks what the
evidence establishes, how the conclusions depend on information and assumptions,
and where another measurement would be more useful than another estimator.
Part IX turns those findings into an investment question: what to build, what
to reuse, and which experiment could change the choice.
